\section{Background}\label{sec:background}

\subsection{What the standards fix and what they leave open}

Under the current expected credit loss standard (CECL, ASC 326), a lender
reserves at each reporting date for the credit losses it expects over the
remaining contractual life of its loans, considering prepayments
\citep{asc326}. The standard is explicit that the method is the lender's
choice: ``The allowance for credit losses may be determined using various
methods'' (ASC 326-20-30-3). Beyond a reasonable and supportable forecast
period the estimate reverts to historical loss information, and the lender may
revert ``immediately, on a straight-line basis, or using another rational and
systematic basis'' (326-20-30-9). The 2023 interagency policy statement
repeats both freedoms and adds that estimating allowances ``involves a high
degree of management judgment and is inherently imprecise''
\citep{ips2023acl}. IFRS 9 differs mainly in staging: lifetime losses are
recognised only once credit risk ``has increased significantly since initial
recognition'' (paragraph 5.5.3), and twelve-month losses otherwise
\citep{ifrs9}.

Three other bodies of rules bear on how much an estimate may move with the
method. The SEC's Staff Accounting Bulletin 119 expects registrants to
``maintain documentation to support the identified range and the rationale
used for determining which estimate is the best estimate within the range of
expected credit losses'' (Question 5; \citealp{sab119}). The bulletin presumes a range exists
but says nothing about how wide it is or where it comes from. For auditors,
AS 2501 asks whether the company has a reasonable basis for selecting its
assumptions from a range of potential assumptions \citep{pcaob2501}, and when
a recorded estimate falls outside a range of reasonable estimates, AS 2810
treats the difference between it and ``the closest reasonable estimate'' as a
misstatement \citep{pcaob2810}. The revised interagency model-risk
guidance of April 2026 (SR 26-2, OCC Bulletin 2026-13), which replaced SR 11-7,
asks for benchmarking against alternative models in validation, but it does
not mention allowances, CECL or accounting estimates; the nearest it comes is a
reference to errors in financial statements in its definition of model risk
\citep{sr262, SR117}. None of these documents says how far equally defensible
model choices can move an allowance. This paper measures it.

\subsection{Related work}

\paragraph{How much CECL estimates vary.} \citet{breeden2017cecl} applies
several CECL method families (vintage, roll-rate, PD$\times$LGD and
discounted-cash-flow approaches) to agency mortgage data and finds lifetime
loss numbers that vary by a factor of two. Those methods are not equally
accurate, so the spread mixes better and worse models.
\citet{breeden2022selection} hold the method family fixed and rank
macroeconomic sensitivity models by information criteria and show that the
choice among well-ranked models moves lifetime PD materially.
\citet{du2018sources} study how modelling choices add variation to allowances
across banks beyond the variation that risk explains, \citet{canals2026cecl} studies model risk
in CECL for consumer credit, and \citet{forrest2024unavoidable} argues that
some model risk in expected-credit-loss models is unavoidable.
\citet{hibbeln2026multiplicity} show predictive multiplicity in recovery-rate
models, the loss-given-default side of the problem. The closest work on the
timing of reserves, \citet{chae2019cecl}, \citet{covas2018cecl} and
\citet{loudis2019cecl}, asks how CECL would have moved allowances and capital
through the 2007--2009 cycle, not how they depend on the model.

\paragraph{Predictive multiplicity.} Many models can fit the same data about
equally well while disagreeing on individual cases: the ``Rashomon effect''
\citep{breiman2001two}. The machine-learning literature measures how often
equally accurate classifiers disagree \citep{marx2020predictive,
hsu2022rashomon, watsondaniels2023prob}, shows that such disagreement comes
from underspecification \citep{damour2022underspec}, and asks what it means for
decisions \citep{black2022multiplicity, semenova2022existence,
fisher2019all}. This paper moves the question from individual decisions
to an aggregate on the balance sheet. Loan-level disagreements could cancel in
a sum over millions of loans, or they could be systematic.

\paragraph{Mortgage default.} Negative equity and adverse shocks drive
mortgage default \citep{foote2008negative, elul2010triggers,
campbell2015model}, and flexible models predict it better than linear ones on
large loan panels \citep{sadhwani2021deep, fuster2022predictably}. The inputs
and the competing-risks structure here follow that literature. Judicial
foreclosure lengthens the time to liquidation
\citep{pence2006foreclosing, gerardi2013borrower}.

\paragraph{Contribution.} I know of no study that builds a formal set of
equally accurate loan-level default models on public data and measures how far
the CECL allowance moves across it, with the data, the economic forecast, the
reversion policy and the loss-severity model held fixed. Doing so at
\num{const.n_dates} year-ends over two decades also shows how the spread
behaves into a crisis, compares it with the discretion the standard grants
explicitly and with the losses that followed, and tests whether it reaches
IFRS 9 staging.
